% Appendix: the proof of parts (i) and (ii) of prop:set-cyclicLR, the large-radius cyclic structure.
\section{The large-radius cyclic structure}
\label{sec:set-lrproof}

This appendix proves parts (i) and (ii) of \cref{prop:set-cyclicLR} from (M), (G0), \cref{lem:set-pf5} and the pairing of \cref{def:set-fq}, in the notation of \cref{sec:set-fq}. Every step is an identity in $\bbQ\fps q$ or $\bbQ\fps q[\logq]$; no truncation in $q$, analytic continuation or conifold frame is used. The letters $\mathfrak w$, $\cP$, $\mathfrak p_k$, $\tilde{\mathfrak p}_m$, $\mathsf p^{(a)}_k$ and $\chi$ are local. Before \cref{prop:set-cyclicLR}~(i) is proved, $\kappa^{\mathrm{LR}}_a$ denotes the structure constant of \cref{lem:set-grading}.

\begin{lem}[The $\bbZ/5$ grading]
\label{lem:set-grading}
At $\lambda_i=\zeta^i$ and $t=TH$, $\langle\!\langle H^a,H^b,H^c\rangle\!\rangle=0$ unless $a+b+c\equiv3\bmod5$. Hence $H^a*H^b$ is a multiple of $H^{a+b}$, and $H*H^a=\kappa^{\mathrm{LR}}_aH^{a+1}$ with $\kappa^{\mathrm{LR}}_a\in1+q\bbQ\fps q$.
\end{lem}

\begin{proof}
By the divisor equation (G0)(b), with $\Qn=1$,
\[
\langle\!\langle H^a,H^b,H^c\rangle\!\rangle(TH)=\eta(H^aH^b,H^c)+\sum_{d\geq1}\re^{dT}\langle H^a,H^b,H^c\rangle_{0,3,d};
\]
in degree $0$ the insertions of $TH$ contribute nothing, as $\overline{\mathcal M}_{0,3+n}(\bbP^4,0)=\overline{\mathcal M}_{0,3+n}\times\bbP^4$ with the integrand pulled back from $\bbP^4$. By (G0)(a), before the specialisation the term of degree $d\geq1$ is a polynomial in the elementary symmetric functions $e_1,\dots,e_5$ of the weights, of weighted degree $m=a+b+c-3$. At $\lambda_j=\zeta^j$, $\prod_j(x-\zeta^j)=x^5-1$ gives $e_1=\dots=e_4=0$ and $e_5=1$, so the value is $0$ unless $5\mid m$. The classical term is $\eta(H^{a+b},H^c)=5\,\delta_{a+b+c\equiv3}$. Since $\eta^{cc'}$ is supported on $c+c'\equiv3$, $H^a*H^b$ is supported on $H^{a+b}$. The structure constant $\kappa^{\mathrm{LR}}_a$ is a power series in $\re^T$ with rational coefficients, hence lies in $\bbQ\fps q$, and $\kappa^{\mathrm{LR}}_a(0)=1$ since at $q=0$ the product is the classical one.
\end{proof}

\begin{lem}[Unit and Frobenius property]
\label{lem:set-unitfrob}
$\kappa^{\mathrm{LR}}_0=\kappa^{\mathrm{LR}}_2=1$ and $\kappa^{\mathrm{LR}}_3=\kappa^{\mathrm{LR}}_4$.
\end{lem}

\begin{proof}
The class $1$ is the unit, so $\kappa^{\mathrm{LR}}_0=1$. By the Frobenius property, $5\kappa^{\mathrm{LR}}_a=\eta(H*H^a,H^{2-a})=\eta(H^a,H*H^{2-a})=5\kappa^{\mathrm{LR}}_{2-a}$ for all $a\in\bbZ/5$, which gives $\kappa^{\mathrm{LR}}_0=\kappa^{\mathrm{LR}}_2$ and $\kappa^{\mathrm{LR}}_3=\kappa^{\mathrm{LR}}_4$.
\end{proof}

\begin{lem}[The periods]
\label{lem:set-periods}
Let $\mathfrak w$ be an indeterminate and
\[
\cP(\mathfrak w,q)\coloneqq\sum_{d\geq0}q^d\,\frac{\prod_{k=1}^{5d}(5\mathfrak w+k)}{\prod_{k=1}^d\bigl((\mathfrak w+k)^5-\mathfrak w^5\bigr)}\in\bbQ\fps q\fps{\mathfrak w},
\]
the denominators being polynomials in $\mathfrak w$ with the constant term $(d!)^5$. The map $\bbQ\fps q[\logq]\fps{\mathfrak w}\to\bbH\otimes\bbQ\fps q[\logq]\fps{z^{-1}}$, $\sum_ka_k\mathfrak w^k\mapsto\sum_ka_kH^kz^{-k}$ (with $H^k\coloneqq H^{k\bmod5}$), is an injective ring homomorphism, and $I/z$ is the image of $\re^{\logq\mathfrak w}\cP(\mathfrak w,q)$. Hence
\[
\frac Iz=\sum_{k\geq0}\mathfrak p_kH^kz^{-k},\qquad\mathfrak p_k\coloneqq\sum_{j=0}^k\frac{\logq^j}{j!}\,\tilde{\mathfrak p}_{k-j},\qquad\tilde{\mathfrak p}_m\coloneqq\coef{\mathfrak w^m}\cP\in\bbQ\fps q,\quad\tilde{\mathfrak p}_m(0)=\delta_{m0},
\]
with $\mathfrak p_0=C_0$ and $\mathfrak p_1=I_1$. Moreover $\PFfour\mathfrak p_k=0$ for $k=0,\dots,3$, and $\PFfour\mathfrak p_4=1$.
\end{lem}

\begin{proof}
The map is multiplicative because $H^5=1$, and injective because distinct powers of $\mathfrak w$ go to distinct powers of $z$; followed by the restriction to $p_i$ it is $G(\mathfrak w)\mapsto G(\zeta^i/z)$. As $\prod_j(\zeta^i-\zeta^j+kz)=(\zeta^i+kz)^5-1$, the powers $z^{-5d}$ cancel in each term $q^d$, and with $q^{\zeta^i/z}=\re^{\logq\zeta^i/z}$ the restriction of the image of $\re^{\logq\mathfrak w}\cP$ to $p_i$ is $I_i/z$; the five restrictions determine the element. Expanding gives $\mathfrak p_k$, and the term $q^0$ of $\cP$ is $1$. The constant term of $\cP$ in $\mathfrak w$ is $C_0$, and its coefficient of $\mathfrak w$ is, by logarithmic differentiation, $\sum_dq^d\frac{(5d)!}{(d!)^5}\bigl(\sum_{k=1}^{5d}\frac5k-\sum_{k=1}^d\frac5k\bigr)=\tilde I_1$; so $\mathfrak p_0=C_0$ and $\mathfrak p_1=I_1$.

\emph{The operator $\PFfive$.} On the term $q^d$ of $I_i$, which is $q^{d+\zeta^i/z}$ times a function of $z$, $z\Dq$ acts by $\zeta^i+dz$, so $(z\Dq)^5-1=\prod_j(z\Dq-\zeta^j)$ multiplies it by $\prod_j(\zeta^i-\zeta^j+dz)$; this cancels the factors $k=d$ of the denominator (and kills the term $d=0$), and leaves $q\prod_{k=1}^5(5\zeta^i+(5d-5+k)z)$ times the term $q^{d-1}$, which is what $q\prod_{k=1}^5(5z\Dq+kz)$ gives on the term $q^{d-1}$. Hence $\PFfive I_i=0$. By \cref{lem:set-pf5} and $I_i/z=\sum_k\zeta^{ik}\mathfrak p_kz^{-k}$,
\[
0=\PFfive(I_i/z)=\sum_{k\geq0}\zeta^{ik}\bigl(z^{5-k}\,\Dq\PFfour\mathfrak p_k-z^{-k}\mathfrak p_k\bigr),
\]
and the coefficients of $z^5,\dots,z^1$ give $\Dq\PFfour\mathfrak p_k=0$ for $k=0,\dots,4$.

\emph{Constants.} The constants of $\Dq$ on $\bbQ\fps q[\logq]$ are $\bbQ$: on $\bbQ\fps q$ this is clear from $\Dq q^n=nq^n$; and if $f=\sum_{j\leq N}a_j\logq^j$ with $a_N\neq0$, $N\geq1$ and $\Dq f=0$, then $a_N\in\bbQ$ and $\Dq a_{N-1}=-Na_N\neq0$, impossible as $\Dq$ maps $\bbQ\fps q$ into $q\bbQ\fps q$. Hence $\PFfour\mathfrak p_k\in\bbQ$ for $k\leq4$, equal to its part of order $q^0$. Taking this part commutes with $\Dq$, which acts on it as $\dd/\dd\logq$, and kills the multiples of $q$; by \eqref{eq:set-PF4} and $\tilde{\mathfrak p}_m(0)=\delta_{m0}$ it maps $\PFfour\mathfrak p_k$ to $(\dd/\dd\logq)^4(\logq^k/k!)$, which is $0$ for $k\leq3$ and $1$ for $k=4$.
\end{proof}

\begin{lem}[The Birkhoff recursion]
\label{lem:set-birkhoff}
Put $\mathsf p^{(0)}_k\coloneqq\mathfrak p_k/C_0$ for $k\geq0$ and, recursively for $a=0,\dots,3$, $\mathsf p^{(a+1)}_k\coloneqq\partial_T\mathsf p^{(a)}_k/\kappa^{\mathrm{LR}}_a$ for $k\geq a+1$. Then, along the line $t=TH$,
\[
J_a=\sum_{k\geq a}\mathsf p^{(a)}_k\,H^kz^{a-k},\qquad\mathsf p^{(a)}_a=1\qquad(a=0,\dots,4),
\]
and $\kappa^{\mathrm{LR}}_a=\partial_T\mathsf p^{(a)}_{a+1}$ for $a=0,\dots,3$. In particular $\mathsf p^{(0)}_1=T$, $\kappa^{\mathrm{LR}}_1=\partial_T^2\mathsf p^{(0)}_2$, $\kappa^{\mathrm{LR}}_2=\partial_T\bigl(\partial_T^2\mathsf p^{(0)}_3/\kappa^{\mathrm{LR}}_1\bigr)$ and, since $\kappa^{\mathrm{LR}}_2=1$, $\kappa^{\mathrm{LR}}_3=\partial_T^2\bigl(\partial_T^2\mathsf p^{(0)}_4/\kappa^{\mathrm{LR}}_1\bigr)$.
\end{lem}

\begin{proof}
Induction on $a$. By (G0)(d), (M) and \cref{lem:set-periods}, $J_0=J/z=I/(zC_0)=\sum_k(\mathfrak p_k/C_0)H^kz^{-k}$ along the line, with $\mathsf p^{(0)}_0=1$; distinct $k$ carry distinct powers of $z$. By (G0)(d) with $b=1$ and \cref{lem:set-grading}, $z\partial_TJ_a=\kappa^{\mathrm{LR}}_aJ_{a+1}$ for $a\leq3$; only derivatives along the line enter. Differentiating the expansion of $J_a$ gives $z\partial_TJ_a=\sum_{k\geq a+1}(\partial_T\mathsf p^{(a)}_k)H^kz^{a+1-k}$, and since $J_{a+1}=H^{a+1}+O(z^{-1})$ the coefficient of $z^0$ gives $\kappa^{\mathrm{LR}}_a=\partial_T\mathsf p^{(a)}_{a+1}$; as $\kappa^{\mathrm{LR}}_a$ is a unit (\cref{lem:set-grading}), $J_{a+1}=z\partial_TJ_a/\kappa^{\mathrm{LR}}_a$ gives the expansion of $J_{a+1}$, with $\mathsf p^{(a+1)}_{a+1}=1$. Finally $\mathsf p^{(0)}_1=I_1/C_0=T$, $\kappa^{\mathrm{LR}}_0=\partial_TT=1$, and the closed forms follow from $\kappa^{\mathrm{LR}}_a=\partial_T\mathsf p^{(a)}_{a+1}$ and \cref{lem:set-unitfrob}.
\end{proof}

\begin{lem}[The Wronskian]
\label{lem:set-wronskian}
$\Wr(\mathfrak p_0,\dots,\mathfrak p_3)\coloneqq\det\bigl(\Dq^j\mathfrak p_k\bigr)_{j,k=0,\dots,3}=L^{10}$.
\end{lem}

\begin{proof}
By \eqref{eq:set-PF4} the leading coefficient of $\PFfour$ is $1-5^5q=L^{-5}$ and its coefficient of $\Dq^3$ is $-6250q$. Differentiating the determinant row by row (the rows $j<3$ give two equal rows) and replacing $\Dq^4\mathfrak p_k$ in the last row by means of $\PFfour\mathfrak p_k=0$ ($k\leq3$, \cref{lem:set-periods}), whose terms of order $\leq2$ are combinations of the other rows, gives $\Dq\Wr=6250qL^5\,\Wr$; also $\Dq L^{10}=10L^9\Dq L=6250qL^5\,L^{10}$. So $\Wr\,L^{-10}$ is a constant of $\Dq$, hence (proof of \cref{lem:set-periods}) equal to its part of order $q^0$. Taking this part is a ring homomorphism to $\bbQ[\logq]$, on which $\Dq$ acts as $\dd/\dd\logq$; it maps $\Dq^j\mathfrak p_k$ to $\logq^{k-j}/(k-j)!$ for $k\geq j$ and to $0$ for $k<j$, so the determinant to $1$.
\end{proof}

\begin{lem}[Operators of order at most three]
\label{lem:set-order3}
If $\Theta=\sum_{j=0}^3r_j\Dq^j$ with $r_j\in\bbQ\lser q$ satisfies $\Theta\mathfrak p_k=0$ for $k=0,\dots,3$, then $\Theta=0$.
\end{lem}

\begin{proof}
As the $r_j$ do not involve $\logq$ and $\Dq\re^{\logq\mathfrak w}=\mathfrak w\re^{\logq\mathfrak w}$, \cref{lem:set-periods} gives, coefficientwise in $\mathfrak w$,
\[
\sum_k(\Theta\mathfrak p_k)\mathfrak w^k=\Theta\bigl(\re^{\logq\mathfrak w}\cP\bigr)=\re^{\logq\mathfrak w}\,\cP_\Theta,\qquad\cP_\Theta\coloneqq\sum_{j=0}^3r_j(\Dq+\mathfrak w)^j\cP\in\bbQ\lser q\fps{\mathfrak w}.
\]
The hypothesis says that the left side is $O(\mathfrak w^4)$, hence so is $\cP_\Theta$. Now $\coef{\mathfrak w^n}\cP_\Theta=\sum_{j=0}^3m_{nj}r_j$ with $m_{nj}\coloneqq\sum_{i\leq\min(j,n)}\binom ji\Dq^{j-i}\tilde{\mathfrak p}_{n-i}\in\bbQ\fps q$, and $m_{nj}\equiv\delta_{nj}\bmod q$ because $\tilde{\mathfrak p}_m(0)=\delta_{m0}$ and $\Dq$ kills $q^0$. So $\det(m_{nj})$ is a unit, and $r_j=0$.
\end{proof}

\begin{proof}[Proof of \cref{prop:set-cyclicLR}~(i) and~(ii)]
(i) By \cref{lem:set-unitfrob}, $\kappa^{\mathrm{LR}}_0=\kappa^{\mathrm{LR}}_2=1$ and $\kappa^{\mathrm{LR}}_3=\kappa^{\mathrm{LR}}_4$.

\emph{The coefficient $\kappa^{\mathrm{LR}}_1$.} By \cref{lem:set-periods,lem:set-birkhoff}, $\mathfrak p_k=C_0\mathsf p^{(0)}_k$ with $\mathsf p^{(0)}_0=1$, $\mathsf p^{(0)}_1=T$. For any four functions $f_k$, $\det(\Dq^j(C_0f_k))=C_0^4\det(\Dq^jf_k)=C_0^4C_1^6\det(\partial_T^jf_k)$, since $\Dq^j(C_0f)-C_0\Dq^jf$ and $\Dq^j-C_1^j\partial_T^j$ are of lower order with coefficients independent of $f$. With $f=(1,T,\mathsf p^{(0)}_2,\mathsf p^{(0)}_3)$ and \cref{lem:set-birkhoff},
\[
\det(\partial_T^jf_k)=\partial_T^2\mathsf p^{(0)}_2\,\partial_T^3\mathsf p^{(0)}_3-\partial_T^3\mathsf p^{(0)}_2\,\partial_T^2\mathsf p^{(0)}_3=\bigl(\partial_T^2\mathsf p^{(0)}_2\bigr)^2\,\partial_T\Bigl(\frac{\partial_T^2\mathsf p^{(0)}_3}{\partial_T^2\mathsf p^{(0)}_2}\Bigr)=(\kappa^{\mathrm{LR}}_1)^2\kappa^{\mathrm{LR}}_2 .
\]
So \cref{lem:set-wronskian} and $\kappa^{\mathrm{LR}}_2=1$ give $(\kappa^{\mathrm{LR}}_1)^2=(L^5/C_0^2C_1^3)^2$ in the domain $\bbQ\fps q$; as $\kappa^{\mathrm{LR}}_1(0)=1=L(0)=C_0(0)=C_1(0)$ (\cref{lem:set-grading}), $\kappa^{\mathrm{LR}}_1=L^5/(C_0^2C_1^3)$.

\emph{The coefficient $\kappa^{\mathrm{LR}}_3$.} The operator $\Mcyc^{\mathrm{LR}}=\partial_T^2\circ(\kappa^{\mathrm{LR}}_1)^{-1}\circ\partial_T^2\circ C_0^{-1}$ has coefficients in $\bbQ\fps q$. By \cref{lem:set-birkhoff}, $\Mcyc^{\mathrm{LR}}\mathfrak p_0=\Mcyc^{\mathrm{LR}}\mathfrak p_1=0$ (as $\partial_T^21=\partial_T^2T=0$), $\Mcyc^{\mathrm{LR}}\mathfrak p_2=\partial_T^2\bigl((\kappa^{\mathrm{LR}}_1)^{-1}\kappa^{\mathrm{LR}}_1\bigr)=0$, $\Mcyc^{\mathrm{LR}}\mathfrak p_3=\partial_T\kappa^{\mathrm{LR}}_2=0$ and $\Mcyc^{\mathrm{LR}}\mathfrak p_4=\kappa^{\mathrm{LR}}_3$. In powers of $\Dq$, $\Mcyc^{\mathrm{LR}}=(\kappa^{\mathrm{LR}}_1C_0C_1^4)^{-1}\Dq^4+(\text{order}\leq3)$, and $\PFfour=L^{-5}\Dq^4+(\text{order}\leq3)$. So $\Mcyc^{\mathrm{LR}}-\chi\,\PFfour$, $\chi\coloneqq L^5/(\kappa^{\mathrm{LR}}_1C_0C_1^4)=C_0/C_1$, has order at most $3$ and kills $\mathfrak p_0,\dots,\mathfrak p_3$ (\cref{lem:set-periods}); by \cref{lem:set-order3} it vanishes. Hence $\kappa^{\mathrm{LR}}_3=\Mcyc^{\mathrm{LR}}\mathfrak p_4=\chi\,\PFfour\mathfrak p_4=\chi=C_0/C_1$, $\kappa^{\mathrm{LR}}_4=\kappa^{\mathrm{LR}}_3$, and
\begin{equation}
\label{eq:set-McycLR}
\Mcyc^{\mathrm{LR}}=\kappa^{\mathrm{LR}}_3\,\PFfour .
\end{equation}

(ii) For a solution $v$, the Frobenius property and the commutativity of $*$ give $z\partial_T\eta(H^a,v)=\eta(H^a,H*v)=\eta(H*H^a,v)=\kappa^{\mathrm{LR}}_a\eta(H^{a+1},v)$, as $\eta(H^a,\cdot)$ has constant coefficients; and $z\partial_T=C_1^{-1}z\Dq$.
\end{proof}

The coefficient $\kappa^{\mathrm{LR}}_3$ sees the relation $H^5=1$, through $\PFfour\mathfrak p_4=1$. Relations equivalent to (i) are stated in \cite[Props.~6.1--6.2]{GR19} and, for $a\leq3$, in \cite[Section~2.4.3, (17), and Section~2.4.4, (19)--(20)]{LP20}; the proof above does not use them.
